\section{Discussion}
\label{sec:disc}

\paragraph{State comparison.} H1 asked whether the S-index variability of HD~166620 in the grand minimum is smaller than in the cycle minima. It is, by both methods and for both minima (Table~\ref{tab:h1}). We treat this as descriptive. It rests on one star, on single intervals, and on an $F$ test whose $p$ values are not credible at face value: the observations inside a season are strongly correlated (minimum~2 has \HOneNMinTwo{} points in four seasons), and the correction subtracts a fixed fraction of the mean as instrumental noise, with a different fraction for the instrument that observed the grand minimum. A lower variance in a lower-activity state is also what one expects without any dynamical transition. The direction of the difference is the finding; its significance is not.

\paragraph{Temporal indicators.} H2 was not supported: for HD~166620 the variance has $\tau \le 0$ for the $\sigma = 1$ residual and for the larger smoothing widths, and the lag-1 autocorrelation has $p \ge 0.05$ in all four configurations. But the test was not possible as designed. The Mount Wilson record of HD~166620 ends in 1995, before the transition, and we could not find the original Mount Wilson observations for 1995--2002 in a public source, so the approach to the grand minimum is not sampled. The other \NSegments{} segments are series with no known transition; they form a null scan, and the absence of a flagged segment is the expected outcome when nothing is there to find. A decisive test needs an annual or finer series that runs through a transition.

\paragraph{Power.} The surrogate test has nominal size, slightly conservative (\AutoUp\,\% at a nominal 5\,\%), and in 192 combinations it flagged two where about ten would be expected by chance if the combinations were independent, which they are not. Its power is low: AUC of \AUCTwentyFour{} to \AUCForty{} for a rising autocorrelation without any cycle, at 24 to 40 years. With a deterministic cycle present, an ARMA model fitted to the series can absorb part of it, so the surrogate null then mimics the cycle and the test is less able to react to changes in it. We did not quantify this. A null outcome here is therefore weak evidence against critical slowing down before a loss of cycle, and the instability of the ARMA fit on short segments (Sect.~\ref{sec:temporal}) adds to that.

\paragraph{What the indicators measure.} A Gaussian smooth of width one year keeps most of a cycle of ten to sixteen years, so the $\sigma = 1$ residual keeps only \ResAmpTen, \ResAmpTwelve{} and \ResAmpSixteen\,\% of the amplitude of cycles of 10, 12 and 16 years. This explains the negative median lag-1 autocorrelation of the residual (\MedACSOne) against \MedACNone{} for the original series. Windows of 12 to 26 points cover about one cycle, so the rolling indicators follow cycle phase more than recovery rate, and the four configurations are not four independent views of the noise. The normalized variance differs between Cycle and Flat stars without any critical slowing down, through the cycle amplitude alone. We attribute the two comparisons that survive within-block Holm correction to this, as an interpretation we did not test; they do not survive Holm over all tests.

\paragraph{Limits.} Segments with interpolated years (up to ten in one star), stars outside the validity range of the $R'_{\rm HK}$ calibration, and the composite character of the solar series are marked in Table~\ref{tab:seg}; we ran no sensitivity check on the interpolation. The Sun is the only segment labeled adequate, but its level can shift where instruments change. Group sizes (four to eight) make every group test low in power, and stars entered through the availability of a series, so the groups are convenience samples. A repeat of the analysis on a series that crosses a known transition, with a power curve at the real segment lengths and with a cycle in the injected series, is the next step.
